% ---------------------------------------------------------------------------------------------------------------
% TEMPLATE PARA TRABALHO DE CONCLUSÃO DE CURSO
% Universidade Tecnológica Federal do Paraná - UTFPR
% Campus Toledo
% Baseado no projeto dos professores Diego Marczal e Michael Vornes
% (http://tcc.tsi.gp.utfpr.edu.br/paginas/modelos-latex-da-utfpr)
% 
% Adaptado Pelo Prof. Felipe W. D. Pfrimer
% Versão: 2018.01
%
%----------------------------------------------------------------------------------------------------------------
% Codificação: UTF-8
% LaTeX:  abnTeX2          
% ---------------------------------------------------------------------------------------------------------------


%CARREGA CLASSE PERSONALIZADA DA UTFPR--------------------------------------------------------------------------
\documentclass[%twoside,                % Impressão em frente e verso
    	        oneside,                % Impressão apenas frente
]{configuracoes/utfpr-abntex2}

% \documentclass[
% 	% -- opções da classe memoir --
% 	12pt,				% tamanho da fonte
% 	openright,			% capítulos começam em pág ímpar (insere página vazia caso preciso)
% 	oneside,			% para impressão em apenas um verso
% 	%twoside,			% para impressão em verso e anverso. Oposto a oneside
% 	a4paper,			% tamanho do papel. 
% 	% -- opções da classe abntex2 --
% 	chapter=TITLE,		% títulos de capítulos convertidos em letras maiúsculas
% 	section=TITLE,		% títulos de seções convertidos em letras maiúsculas
% 	subsection=TITLE,	% títulos de subseções convertidos em letras maiúsculas
% 	subsubsection=TITLE,% títulos de subsubseções convertidos em letras maiúsculas
% 	% -- opções do pacote babel --
% 	english,			% idioma adicional para hifenização
% 	french,				% idioma adicional para hifenização
% 	spanish,			% idioma adicional para hifenização
% 	brazil				% o último idioma é o principal do documento
% 	]{configuracoes/utfpr-abntex2}

% DEFINE A NATUREZA DO TRABALHO (TCC1 OU TCC2)-------------------------------------------------------------------
\newif\iftcctype
\tcctypetrue  % Comente essa linha para TCC2

% INCLUI ARQUIVOS DE CONFIGURAÇÕES-------------------------------------------------------------------------------
% REFERÊNCIAS------------------------------------------------------------------
\usepackage[%
    alf,
    abnt-emphasize=bf,
    bibjustif,
    recuo=0cm,
    abnt-url-package=url,       % Utiliza o pacote url
    abnt-refinfo=yes,           % Utiliza o estilo bibliográfico abnt-refinfo
    abnt-etal-cite=3,
    abnt-etal-list=3,
    abnt-thesis-year=final
]{abntex2cite}                  % Configura as citações bibliográficas conforme a norma ABNT

% PACOTES----------------------------------------------------------------------
\usepackage[bottom]{footmisc}                               % Mantém as notas de rodapé sempre na mesma posição
\usepackage[utf8]{inputenc}                                 % Codificação do documento
\usepackage[T1]{fontenc}                                    % Seleção de código de fonte
\usepackage{booktabs}                                       % Réguas horizontais em tabelas
\usepackage{color, colortbl}                                % Controle das cores
\usepackage{float}                                          % Necessário para tabelas/figuras em ambiente multi-colunas
\usepackage{graphicx}                                       % Inclusão de gráficos e figuras
\usepackage{icomma}                                         % Uso de vírgulas em expressões matemáticas
\usepackage{indentfirst}                                    % Indenta o primeiro parágrafo de cada seção
\usepackage{microtype}                                      % Melhora a justificação do documento
\usepackage{multirow, array}                                % Permite tabelas com múltiplas linhas e colunas
\usepackage{subeqnarray}                                    % Permite subnumeração de equações
\usepackage{lastpage}                                       % Para encontrar última página do documento
\usepackage{verbatim}                                       % Permite apresentar texto tal como escrito no documento, ainda que sejam comandos Latex
\usepackage{amsfonts, amssymb, amsmath}                     % Fontes e símbolos matemáticos
\usepackage[algoruled, portuguese]{algorithm2e}             % Permite escrever algoritmos em português
\usepackage{ae, aecompl}                                    % Fontes de alta qualidade
\usepackage{latexsym}                                       % Símbolos matemáticos
\usepackage{lscape}                                         % Permite páginas em modo "paisagem"
\usepackage{hyperref}
%\usepackage{subfig}                                        % Posicionamento de figuras

% Fontes
%\usepackage[scaled]{helvet}                                 % Usa a fonte Helvetica
\usepackage{times}                                         % Usa a fonte Times
%\usepackage{palatino}                                      % Usa a fonte Palatino
%\usepackage{lmodern}                                       % Usa a fonte Latin Modern
%\usepackage{picinpar}                                      % Dispor imagens em parágrafos
%\usepackage{scalefnt}                                      % Permite redimensionar tamanho da fonte
%\usepackage{upgreek}                                       % Fonte letras gregas

% Redefine a fonte para uma fonte similar a Arial (fonte Helvetica)
%\renewcommand*\familydefault{\sfdefault}

% Para a lista de siglas
\usepackage{nomencl} % Pacote necessário para a lista de siglas.
% O comando a seguir diz ao Latex para salvar as siglas em um arquivo separado:
\fazlistasiglas[\listadesiglasname] % Parâmetro opcional é o título da lista.

% Adicione novos pacotes e macros aqui
\newcommand{\numero}{N\textsuperscript{\underline{o}} }
%\usepackage{listings}

% CONFIGURAÇÕES DE APARÊNCIA DO PDF FINAL--------------------------------------
\makeatletter
\hypersetup{%
    portuguese,
    colorlinks=true,   % true: "links" coloridos; false: "links" em caixas de texto
    linkcolor=black,    % Define cor dos "links" internos
    citecolor=black,    % Define cor dos "links" para as referências bibliográficas
    filecolor=black,    % Define cor dos "links" para arquivos
    urlcolor=black,     % Define a cor dos "hiperlinks"
    breaklinks=true,
    pdftitle={\@title},
    pdfauthor={\@author},
    pdfkeywords={abnt, latex, abntex, abntex2}
}
\makeatother

% ALTERA O ASPECTO DA COR AZUL--------------------------------------------------
\definecolor{blue}{RGB}{41,5,195}

% REDEFINIÇÃO DE LABELS---------------------------------------------------------
\renewcommand{\algorithmautorefname}{Algoritmo}
\def\equationautorefname~#1\null{Equa\c c\~ao~(#1)\null}

% CRIA ÍNDICE REMISSIVO---------------------------------------------------------
\makeindex

% HIFENIZAÇÃO DE PALAVRAS QUE NÃO ESTÃO NO DICIONÁRIO---------------------------
\hyphenation{%
    qua-dros-cha-ve
    Kat-sa-gge-los
}


% INCLUI ARQUIVOS DO TRABALHO DE CONCLUSÃO DE CURSO (PRÉ-TEXTUAIS, TEXTUAIS, PÓS-TEXTUAIS)-----------------------

% INSERE CAPA E FOLHA DE ROSTO
% CAPA---------------------------------------------------------------------------------------------------

% ORIENTAÇÕES GERAIS-------------------------------------------------------------------------------------
% Caso algum dos campos não se aplique ao seu trabalho, como por exemplo,
% se não houve coorientador, apenas deixe vazio.
% Exemplos: 
% \coorientador{}
% \departamento{}

% DADOS DO TRABALHO--------------------------------------------------------------------------------------
\titulo{Título do Trabalho: Subtítulo do Trabalho}
\titleabstract{Title in English}	% Apenas para TCC2
\autor{Nome Completo do Autor}
\autorcitacao{SOBRENOME, Nome} % Sobrenome em maiúsculo
\local{Toledo}
\data{\the\year}

% NATUREZA DO TRABALHO (não alterar)---------------------------------------------------------------------
% Opções: 
% - Trabalho de Conclusão de Curso (se for TCC2)
% - Projeto de Trabalho de Conclusão de Curso (se for TCC1)
\iftcctype
	% TCC1:
	\projeto{Projeto de Trabalho de Conclusão de Curso}
\else
	% TCC2:
	\projeto{Trabalho de Conclusão de Curso}
\fi

% TÍTULO ACADÊMICO---------------------------------------------------------------------------------------
% Opções:
% - Bacharel ou Tecnólogo (Se a natureza for Trabalho de Conclusão de Curso)
\tituloAcademico{Bacharel em Engenharia Eletrônica}

% ÁREA DE CONCENTRAÇÃO E LINHA DE PESQUISA---------------------------------------------------------------
% Se a natureza for Trabalho de Conclusão de Curso, deixe ambos os campos vazios
% Se for programa de Pós-graduação, indique a área de concentração e a linha de pesquisa
\areaconcentracao{}
\linhapesquisa{}

% DADOS DA INSTITUIÇÃO-----------------------------------------------------------------------------------
% Se a natureza for Trabalho de Conclusão de Curso, coloque o nome do curso de graduação em "programa"
% Formato para o logo da Instituição: \logoinstituicao{<escala>}{<caminho/nome do arquivo>}
\instituicao{Universidade Tecnológica Federal do Paraná}
%\departamento{Coordenação de Engenharia Eletrônica}
\programa{Curso de Engenharia Eletrônica}
\logoinstituicao{0.2}{dados/figuras/logo_UTFPR_cor.eps} 

% DADOS DOS ORIENTADORES---------------------------------------------------------------------------------
\orientador{Nome do Orientador}
%\orientador[Orientadora:]{Nome da orientadora}
\instOrientador{{\imprimirinstituicao}}

\coorientador{}
%\coorientador[Coorientadora:]{Nome da coorientadora}
\instCoorientador{}

% FOLHA DE ROSTO--------------------------------------------------------------------------------------------------------

% TRABALHO DE CONCLUSÃO DE CURSO
 \iftcctype
 \preambulo{{\imprimirprojeto} apresentado à disciplina de Trabalho de Conclusão de Curso 1 do {\imprimirprograma} da {\imprimirinstituicao} - UTFPR Campus Toledo, como requisito parcial para a obtenção do título de {\imprimirtituloAcademico}.} 
 \else
 \preambulo{{\imprimirprojeto} apresentado ao {\imprimirprograma} da {\imprimirinstituicao} - UTFPR Campus Toledo, como requisito parcial para a obtenção do título de {\imprimirtituloAcademico}.}
 \fi
% DISSERTAÇÃO DE MESTRADO
% \preambulo{{\imprimirprojeto} apresentada ao Programa de \mbox{Pós-graduação} da {\imprimirinstituicao}, como requisito parcial para obtenção do título de {\imprimirtituloAcademico}.}

% TESE DE DOUTORADO
% \preambulo{{\imprimirprojeto} apresentada ao Programa de \mbox{Pós-graduação} da {\imprimirinstituicao}, como requisito parcial para a obtenção do título de {\imprimirtituloAcademico}.}

% PROJETO DE QUALIFICAÇÃO DE MESTRADO OU DOUTORADO
%\preambulo{{\imprimirprojeto} apresentado ao Programa de \mbox{Pós-graduação} da {\imprimirinstituicao}, como requisito parcial para a obtenção do título de {\imprimirtituloAcademico}.}

% OBSERVAÇÕES-----------------------------------------------------------------------------------------------------------
% Altere este arquivo APENAS comentando as linhas que não se aplicam ao tipo de trabalho acadêmico desejado.


\begin{document}

\pretextual
\imprimircapa    						% Comando para imprimir Capa
\imprimirfolhaderosto{}             	% Comando para imprimir Folha de rosto
%% Isto é um exemplo de Ficha Catalográfica, ou ``Dados internacionais de
% catalogação-na-publicação''. Você pode utilizar este modelo como referência. 
% Porém, provavelmente a biblioteca da sua universidade lhe fornecerá um PDF
% com a ficha catalográfica definitiva após a defesa do trabalho. Quando estiver
% com o documento, salve-o como PDF no diretório do seu projeto e substitua todo
% o conteúdo de implementação deste arquivo pelo comando abaixo:
%
% \begin{fichacatalografica}
%     \includepdf{fig_ficha_catalografica.pdf}
% \end{fichacatalografica}
\begin{fichacatalografica}
	\vspace*{\fill}					% Posição vertical
	\hrule							% Linha horizontal
	\begin{center}					% Minipage Centralizado
	\begin{minipage}[c]{12.5cm}		% Largura
	
	\imprimirautor
	
	\hspace{0.5cm} \imprimirtitulo  / \imprimirautor. --
	\imprimirlocal, \imprimirdata-
	
	\hspace{0.5cm} \pageref{LastPage} p. : il. (algumas color.) ; 30 cm.\\
	
	\hspace{0.5cm} \imprimirorientadorRotulo~\imprimirorientador\\
	
	\hspace{0.5cm}
	\parbox[t]{\textwidth}{\imprimirtipotrabalho~--~\imprimirinstituicao,
	\imprimirdata.}\\
	
	\hspace{0.5cm}
		1. Palavra-chave1.
		2. Palavra-chave2.
		I. Orientador.
		II. Universidade xxx.
		III. Faculdade de xxx.
		IV. Título\\ 			
	
	\hspace{8.75cm} CDU 02:141:005.7\\
	
	\end{minipage}
	\end{center}
	\hrule
\end{fichacatalografica}     % Ficha catalográfica (opcional)
% Obs: a capa e folha de rosto são gerados de acordo com \tcctype

\iftcctype % Se TCC1
% INSERE ELEMENTOS TEXTUAIS DO TCC1 (os elementos opcionais podem ser comentados)

% RESUMO--------------------------------------------------------------------------------

\begin{resumo}[RESUMO]
\begin{SingleSpacing}

% Não altere esta seção do texto--------------------------------------------------------
%\imprimirautorcitacao. \imprimirtitulo. \imprimirdata. \pageref {LastPage} f. \imprimirprojeto\ – \imprimirprograma, \imprimirinstituicao. \imprimirlocal, \imprimirdata.\\
%---------------------------------------------------------------------------------------

O Resumo é um elemento obrigatório em tese, dissertação, monografia e TCC, constituído de uma seqüência de frases concisas e objetivas, fornecendo uma visão rápida e clara do conteúdo do estudo. O texto deverá conter no máximo 500 palavras e ser antecedido pela referência do estudo. Também, não deve conter citações. O resumo deve ser redigido em parágrafo único, espaçamento simples e seguido das palavras representativas do conteúdo do estudo, isto é, palavras-chave, em número de três a cinco, separadas entre si por ponto e finalizadas também por ponto. Usar o verbo na terceira pessoa do singular, com linguagem impessoal, bem como fazer uso, preferencialmente, da voz ativa. Texto contendo um único parágrafo.\\

\textbf{Palavras-chave}: Palavra. Segunda Palavra. Outra palavra.

\end{SingleSpacing}
\end{resumo}

% OBSERVAÇÕES---------------------------------------------------------------------------
% Altere o texto inserindo o Resumo do seu trabalho.
% Escolha de 3 a 5 palavras ou termos que descrevam bem o seu trabalho 
										% Resumo do TCC1
% Lista de Figuras----------------------------------------------------------------

\pdfbookmark[0]{\listfigurename}{lof}
\listoffigures*
\cleardoublepage

% OBSERVAÇÕES---------------------------------------------------------------------
% Este arquivo não precisa de ser alterado, pois a lista é gerada automaticamente.
   	% Lista de Figuras
% LISTA DE QUADROS----------------------------------------------------------------

\renewcommand{\listofquadrosname}{LISTA DE QUADROS}

\pdfbookmark[0]{\listofquadrosname}{loq}
\listofquadros*
\cleardoublepage

% OBSERVAÇÕES---------------------------------------------------------------------
% Este arquivo não necessita de ser editado. A lista é gerada automaticamente.
   	% Lista de Quadros (opcional)
% LISTA DE TABELAS-------------------------------------------------------------

\pdfbookmark[0]{\listtablename}{lot}
\listoftables*
\cleardoublepage

% OBSERVAÇÕES-------------------------------------------------------------------
% Este arquivo não precisa ser alterado, pois a lista é gerada automaticamente.
    % Lista de Tabelas
% LISTA DE ABREVIATURAS E SIGLAS----------------------------------------------------------

%\begin{siglas}
%    \item[ABNT] Associação Brasileira de Normas Técnicas
%    \item[DECOM] Departamento de Computação
%\end{siglas}

% OBSERVAÇÕES-----------------------------------------------------------------------------
% Altere a lista acima para definir os acrônimos e siglas utilizados neste trabalho


\listadesiglas     % Lista de Abreviaturas e Siglas
% LISTA DE SÍMBOLOS------------------------------------------------------------

%\begin{simbolos}
%    \item[$ \Gamma $] Letra grega Gama
%    \item[$ \lambda $] Comprimento de onda
%    \item[$ \in $] Pertence
%\end{simbolos}

% OBSERVAÇÕES-------------------------------------------------------------------
% Altere a lista acima para definir os símbolos utilizados no trabalho
\listadesimbolos   % Lista de Símbolos
% LISTA DE ALGORITMOS----------------------------------------------------------

\newcommand{\algoritmoname}{Algoritmo}
\renewcommand{\listalgorithmcfname}{LISTA DE ALGORITMOS}

\floatname{algocf}{\algoritmoname}
\newlistof{listofalgoritmos}{loa}{\listalgoritmoname}
\newlistentry{algocf}{loa}{0}

\counterwithout{algocf}{chapter}
\renewcommand{\cftalgocfname}{\algoritmoname\space}
\renewcommand*{\cftalgocfaftersnum}{\hfill--\hfill}

\pdfbookmark[0]{\listalgorithmcfname}{loa}
\listofalgorithms
\cleardoublepage

% OBSERVAÇÕES------------------------------------------------------------------
% Este arquivo não precisa ser alterado, pois a lista é gerada automaticamente.
 % Lista de Algoritmos (opcional)
% SUMÁRIO----------------------------------------------------------------------

\renewcommand{\contentsname}{SUMÁRIO}

\pdfbookmark[0]{\contentsname}{toc}
\tableofcontents*
\cleardoublepage

% OBSERVAÇÕES-------------------------------------------------------------------
% Este arquivo não precisa ser alterado, pois o sumário é gerado automaticamente.
               	% Sumário
%---------texto-----------------------------------------------------------------
\textual
% INTRODUÇÃO-------------------------------------------------------------------

\chapter{INTRODUÇÃO}
\label{chap:introducao}

Edite e coloque aqui o seu texto de introdução.

A Introdução é a parte inicial do texto, na qual devem constar o tema e a delimitação do assunto tratado, objetivos da pesquisa e outros elementos necessários para situar o tema do trabalho, tais como: justificativa, procedimentos metodológicos (classificação inicial), embasamento teórico (principais bases sintetizadas) e estrutura do trabalho, tratados de forma sucinta. Recursos utilizados e cronograma são incluídos quando necessário. Salienta-se que os procedimentos metodológicos e o embasamento teórico são tratados, posteriormente, em capítulos próprios e com a profundidade necessária ao trabalho de pesquisa.

\section{LEIA ESTA SEÇÃO ANTES DE COMEÇAR}
\label{sec:antesleiame}

Este documento é um \emph{template} \LaTeX{} que foi concebido, primariamente, para ser utilizado na elaboração de Trabalho de Conclusão de Curso \sigla[Trabalho de Conclusão de Curso]{TCC} em conformidade com as normas da Associação Brasileira de Normas técnicas \sigla[Associação Brasileira de Normas técnicas]{ABNT}.

Para a produção deste \emph{template} foi necessário adaptar o arquivo {\ttfamily abntex2.cls}. Assim, foi produzido o arquivo {\ttfamily utfpr-abntex2.cls} que define o \verb|documentclass| específico para a Universidade Tecnológica Federal do Paraná \sigla[Universidade Tecnológica Federal do Paraná]{UTFPR}.

Antes de começar a escrever o seu trabalho acadêmico utilizando este \emph{template}, é importante saber que há dois arquivos que você precisará editar para que a capa e a folha de rosto de seu trabalho sejam geradas automaticamente.
São eles os arquivos {\ttfamily utfpr-tcc.tex} e {\ttfamily capa.tex}, presentes no diretório raiz e {\ttfamily /comum}, respectivamente. 

No arquivo {\ttfamily utfpr-tcc.tex}, que contém toda a estrutura do documento, deverá ser informado a natureza do trabalho (TCC1 ou TCC2). Para isso, procure o seguinte trecho:

\begin{verbatim}	
		\newif\iftcctype
		\tcctypetrue	% Comente essa linha para TCC2		
\end{verbatim}

\noindent Comente a segunda linha para habilitar a compilação dos arquivos do TCC2, e não comente para os arquivos de TCC1. 

No arquivo {\ttfamily capa.tex} deverá ser informado nome do autor, título do trabalho, nome do orientador e outras informações necessárias.

Em sistemas Linux, a compilação para gerar um arquivo no formato pdf, incluindo corretamente as referências bibliográficas, deve ser realizada utilizando o comando \verb|makefile|, disponível na mesma pasta onde está o arquivo principal \verb|utfpr-tcc.tex|. Caso seja alterado o nome do arquivo \verb|utfpr-tcc.tex|, deverá ser alterado no arquivo \verb|makefile| também.

\section{ORGANIZAÇÃO DO TRABALHO}
\label{sec:organizacaoTrabalho}

Normalmente ao final da introdução é apresentada, em um ou dois parágrafos curtos, a organização do restante do trabalho acadêmico.
Deve-se dizer o quê será apresentado em cada um dos demais capítulos.

%\simb[texto]{x_i}  

%\sigla[exemplo]{EX}
				% Introdução
% DELIMITAÇÃO DO TEMA-------------------------------------------------------------------

\chapter{DELIMITAÇÃO DO TEMA}
\label{chap:delimitacao}

A delimitação do tema é um capítulo opcional, normalmente exigido em editais de bolsa para TCC da UTFPR. Para mais informações acesse a página da web: \href{https://guiadamonografia.com.br/como-fazer-a-delimitacao-do-tema/}{O Guia da Monografia}.
				% Delimitação do Tema (opcional - edital de tcc)
% PROBLEMAS E PREMISSAS-------------------------------------------------------------------

\chapter{PROBLEMAS E PREMISSAS}
\label{chap:problemas}

Este capítulo deve tratar da dificuldade teórica ou prática a ser resolvida pelo trabalho. Inclui Hipóteses ou Premissas (Pressuposições \/Conjecturas que antecedem a constatação dos fatos). Também é um capítulo opcional, normalmente exigido em editais de bolsa para TCC da UTFPR.
				% Problemas e Premissas (opcional - edital de tcc)
% OBJETIVOS------------------------------------------------------------------

\chapter{OBJETIVOS}
\label{chap:objetivos}

Os objetivos devem deixar claro o que se pretende alcançar com o trabalho proposto. 

\section{OBJETIVO GERAL}
\label{sec:obgeral}

Deve salientar o objetivo mais amplo do trabalho. Trata-se de uma ação e, portanto, deve ser construído sempre a partir de um verbo no infinitivo.

\section{OBJETIVOS ESPECÍFICOS}
\label{sec:obesp}

São as etapas a serem realizadas para que se atinja o objetivo geral. Sugere-se que na descrição dos objetivos específicos sejam utilizados verbos de ação:

\begin{itemize}
    \item Objetivos Específicos Descritivos: caracterizar, descrever, traçar, ...;
    \item Objetivos Específicos Exploratórios: conhecer, identificar, levantar, ...;   
    \item Objetivos Específicos Explicativos: analisar, avaliar, verificar, ...
\end{itemize}
                % Objetivos
% JUSTIFICATIVA------------------------------------------------------------------

\chapter{JUSTIFICATIVA}
\label{chap:justificativa}

Na justificativa deve estar claro:

\begin{itemize}
    \item Por que se pretende resolver o problema apresentado;
    \item Qual a relevância do trabalho dentro da área técnica;   
    \item Quais os possíveis ganhos que se pode agregar ao conhecimento técnico com a resolução deste problema;
		\item Quais os fatores motivadores para se dedicar a este problema específico.
\end{itemize}
            % Justificativa
% REVISÃO DE LITERATURA--------------------------------------------------------

\chapter{REVISÃO DA LITERATURA}
\label{chap:fundamentacaoTeorica}

Uma dica: é uma boa prática iniciar cada novo capítulo com um breve texto introdutório (tipicamente, dois ou três parágrafos) que deve deixar claro o quê será discutido no capítulo, bem como a organização do capítulo.

Este capítulo também pode ser chamado de ``REFERENCIAL TEÓRCIO''. Para mais informações sobre como estruturar um referencial teórico e evitar plágios, acesse a página: \href{https://guiadamonografia.com.br/como-fazer-referencial-teorico/}{O Guia da Monografia}


      % Referencial teorico
% MATERIAIS E MÉTODOS------------------------------------------------------------------

\chapter{MATERIAIS E MÉTODOS}
\label{chap:materiaisemetodos}

Este capítulo deve tratar de como se pretende realizar o trabalho/pesquisa:

\begin{itemize}
    \item Etapas;
    \item Procedimentos de estudos e coleta de dados;   
    \item Estratégias a serem utilizadas para análise de dados, realização de ensaios, etc.;
		\item Materiais que serão utilizados.
\end{itemize}        % Materiais e metodologia
% DEMONSTRATIVO DE APLICAÇÃO DE RECURSOS-------------------------------------------------------------------

\chapter{DEMONSTRATIVO DE APLICAÇÃO DE RECURSOS}
\label{chap:demonstrativo}

O demonstrativo de aplicação de recursos é um capítulo opcional, normalmente exigido em editais de bolsa para TCC da UTFPR.


           % Demonstrativo de recursos (opcional - edital de tcc)
% RESULTADOS ESPERADOS-------------------------------------------------------------------

\chapter{RESULTADOS ESPERADOS}
\label{resultadosesperados}

É nesse capítulo onde são colocados as expectativas do projeto. Para maiores informações acesse a página: \href{http://www.mundodamonografia.com.br/projeto-de-pesquisa-resultados/}{Mundo da Monografia}.

        		% Resultados esperados
% CRONOGRAMA-------------------------------------------------------------------

\chapter{CRONOGRAMA}
\label{chap:cronograma}

Este capítulo detalha que etapas serão realizadas ao longo do tempo. O cronograma se refere ao período de um ano: TCC1 e TCC2, onde todas as atividades a serem desenvolvidas devem ser citadas e detalhadas. Segue um exemplo de como o cronograma deve ser estruturado:

% Insere uma linha vertical
\noindent\hrulefill

O desenvolvimento deste trabalho se dará da seguinte forma:

% Atividades
\begin{enumerate}
	\item \label{ela-pro} \textbf{Elaboração da proposta de TCC}: descreva aqui a atividade.
	\item \label{anI} \textbf{Análise dos métodos de localização de regiões candidatas em imagens digitais e
		reconhecimento de caracteres através de redes neurais}: descreva aqui a atividade.
	\item \label{anII} \textbf{Análise do funcionamento da técnica em microcontroladores}: descreva aqui a atividade.
	\item \label{anIII} \textbf{Análise dos modelos implementados em hardware}: descreva aqui a atividade. 
	\item \label{dI} \textbf{Estudo do Linux Embarcado}: descreva aqui a atividade.
	\item \label{dII}  \textbf{Validação dos módulos de localização de regiões candidatas}: descreva aqui a atividade..
	\item \label{dIII} \textbf{Implementação do módulo de captura de imagem}: descreva aqui a atividade. 
	\item \label{esc-tcI}  \textbf{Escrita do TCC1}: descreva aqui a atividade.
	\item \label{imI} \textbf{Implementação da máquina de estados para controle da análise de imagem}: descreva aqui a atividade.
	\item \label{imII} \textbf{Desenvolvimento da camada de integração}: descreva aqui a atividade.
	\item \label{imIII} \textbf{Integração dos módulos que compõe o sistema}: descreva aqui a atividade.
	\item \label{tec} \textbf{Teste e correções}: descreva aqui a atividade.
	\item \label{esc-tcII} \textbf{Escrita do TCC2}: descreva aqui a atividade.
\end{enumerate}

O \autoref{qua:cronograma} mostra o período previsto para as atividades propostas.

% insere o cronograma
% CRONOGRAMA-------------------------------------------------------------------

\chapter{CRONOGRAMA}
\label{chap:cronograma}

Este capítulo detalha que etapas serão realizadas ao longo do tempo. O cronograma se refere ao período de um ano: TCC1 e TCC2, onde todas as atividades a serem desenvolvidas devem ser citadas e detalhadas. Segue um exemplo de como o cronograma deve ser estruturado:

% Insere uma linha vertical
\noindent\hrulefill

O desenvolvimento deste trabalho se dará da seguinte forma:

% Atividades
\begin{enumerate}
	\item \label{ela-pro} \textbf{Elaboração da proposta de TCC}: descreva aqui a atividade.
	\item \label{anI} \textbf{Análise dos métodos de localização de regiões candidatas em imagens digitais e
		reconhecimento de caracteres através de redes neurais}: descreva aqui a atividade.
	\item \label{anII} \textbf{Análise do funcionamento da técnica em microcontroladores}: descreva aqui a atividade.
	\item \label{anIII} \textbf{Análise dos modelos implementados em hardware}: descreva aqui a atividade. 
	\item \label{dI} \textbf{Estudo do Linux Embarcado}: descreva aqui a atividade.
	\item \label{dII}  \textbf{Validação dos módulos de localização de regiões candidatas}: descreva aqui a atividade..
	\item \label{dIII} \textbf{Implementação do módulo de captura de imagem}: descreva aqui a atividade. 
	\item \label{esc-tcI}  \textbf{Escrita do TCC1}: descreva aqui a atividade.
	\item \label{imI} \textbf{Implementação da máquina de estados para controle da análise de imagem}: descreva aqui a atividade.
	\item \label{imII} \textbf{Desenvolvimento da camada de integração}: descreva aqui a atividade.
	\item \label{imIII} \textbf{Integração dos módulos que compõe o sistema}: descreva aqui a atividade.
	\item \label{tec} \textbf{Teste e correções}: descreva aqui a atividade.
	\item \label{esc-tcII} \textbf{Escrita do TCC2}: descreva aqui a atividade.
\end{enumerate}

O \autoref{qua:cronograma} mostra o período previsto para as atividades propostas.

% insere o cronograma
% CRONOGRAMA-------------------------------------------------------------------

\chapter{CRONOGRAMA}
\label{chap:cronograma}

Este capítulo detalha que etapas serão realizadas ao longo do tempo. O cronograma se refere ao período de um ano: TCC1 e TCC2, onde todas as atividades a serem desenvolvidas devem ser citadas e detalhadas. Segue um exemplo de como o cronograma deve ser estruturado:

% Insere uma linha vertical
\noindent\hrulefill

O desenvolvimento deste trabalho se dará da seguinte forma:

% Atividades
\begin{enumerate}
	\item \label{ela-pro} \textbf{Elaboração da proposta de TCC}: descreva aqui a atividade.
	\item \label{anI} \textbf{Análise dos métodos de localização de regiões candidatas em imagens digitais e
		reconhecimento de caracteres através de redes neurais}: descreva aqui a atividade.
	\item \label{anII} \textbf{Análise do funcionamento da técnica em microcontroladores}: descreva aqui a atividade.
	\item \label{anIII} \textbf{Análise dos modelos implementados em hardware}: descreva aqui a atividade. 
	\item \label{dI} \textbf{Estudo do Linux Embarcado}: descreva aqui a atividade.
	\item \label{dII}  \textbf{Validação dos módulos de localização de regiões candidatas}: descreva aqui a atividade..
	\item \label{dIII} \textbf{Implementação do módulo de captura de imagem}: descreva aqui a atividade. 
	\item \label{esc-tcI}  \textbf{Escrita do TCC1}: descreva aqui a atividade.
	\item \label{imI} \textbf{Implementação da máquina de estados para controle da análise de imagem}: descreva aqui a atividade.
	\item \label{imII} \textbf{Desenvolvimento da camada de integração}: descreva aqui a atividade.
	\item \label{imIII} \textbf{Integração dos módulos que compõe o sistema}: descreva aqui a atividade.
	\item \label{tec} \textbf{Teste e correções}: descreva aqui a atividade.
	\item \label{esc-tcII} \textbf{Escrita do TCC2}: descreva aqui a atividade.
\end{enumerate}

O \autoref{qua:cronograma} mostra o período previsto para as atividades propostas.

% insere o cronograma
\input{dados/quadros/cronograma/cronograma.tex}





        		% Resultados esperados
% ORIENTAÇÕES GERAIS------------------------------------------------------------
\chapter{ORIENTAÇÕES GERAIS SOBRE O MODELO}
\label{chap:orientações}

Este capítulo trata de orientações gerais sobre como utilizar este modelo de TCC. Recomenda-se a leitura do texto compilado junto com o código \LaTeX{} deste capítulo.

% SOBRE AS LISTAS DE SÍMBOLOS E SIGLAS------------------------------------------
\section{SOBRE AS LISTAS DE SÍMBOLOS E SIGLAS} 
\label{sec:simbsiglas}

Símbolos são inseridos automaticamente na lista através do comando:

\begin{verbatim}
\simb[<significado>]{<símbolo>}
\end{verbatim}

%Exemplos de código são mostrados no \autoref{qua:simb}.

%{\input{./dados/quadros/simb.tex}}

Dessa forma, o código

\begin{verbatim}
\simb[Variação da tensão de saída]{\Delta V_x}
\end{verbatim}

\noindent resulta em ``$\Delta V_x$'' no texto e a adição deste termo na lista de símbolos.

Siglas e abreviaturas são inseridos automaticamente na lista através do comando:

\begin{verbatim}
\sigla[<significado>]{<sigla ou abreviatura>}
\end{verbatim}

Dessa forma, o código

\begin{verbatim}
\simb[Corrente Alternada]{CA}
\end{verbatim}

\noindent resulta em ``(CA)'' no texto e a adição deste termo na lista de Abreviaturas e Siglas. Note que a silga ou abreveatura fica entre parenteses.


% SOBRE AS REFERENCIAS CRUZADAS E ITENS NUMERADOS-------------------------------
\section{SOBRE AS REFERENCIAS CRUZADAS E ITENS NUMERADOS}
\label{sec:refs}

A numeração sequencial de figuras, tabelas, quadros, algoritmos e equações ocorre de modo automático. As ilustrações (figuras), tabelas, algoritmos e quadros serão indexadas automaticamente em suas respectivas listas. 

Referências cruzadas são obtidas através dos comandos \verb|\label{}| e \verb|\ref{}|. Sendo assim, não é necessário por exemplo, saber que o número de certo capítulo é \ref{chap:fundamentacaoTeorica} para colocar o seu número no texto. Outra forma que pode ser utilizada é esta: \autoref{chap:fundamentacaoTeorica}, facilitando a inserção, remoção e manejo de elementos numerados no texto sem a necessidade de renumerar todos esses elementos.

% FIGURAS-----------------------------------------------------------------------
\section{FIGURAS}
\label{sec:figuras}

Esta seção trata de um exemplo de como inserir uma figura. Note que a \autoref{fig:figura-exemplo1} aparece automaticamente na lista de figuras. Para saber mais sobre o uso de imagens no \LaTeX{} consulte literatura especializada \cite{Goossens2007}.

Os arquivos das figuras devem ser armazenados no diretório de "/dados", para facilitar a organização.

\begin{figure}[!htb]
    \centering
    \caption{Exemplo de Figura}
    \includegraphics[width=0.6\textwidth]{dados/figuras/orientacoes/blackmesalogo.eps}
    \fonte{\citeonline{valve1998}}
    \label{fig:figura-exemplo1}
\end{figure}

% QUADROS E TABELAS---------------------------------------------------------------
\section{QUADROS E TABELAS}
\label{sec:tabelas}

Esta seção trata de exemplos de como inserir o \autoref{qua:quadro-exemplo1} e a \autoref{tab:tabela-exemplo1}. Ambos aparecem automaticamente nas suas respectivas listas. Para saber mais informações sobre a construção de tabelas no \LaTeX{} consulte literatura especializada \cite{Mittelbach2004}.

Ambos os elementos (Quadros e Tabelas) devem ser criados em arquivos separados para facilitar manutenção e armazenados no diretório de "/dados".

\begin{quadro}[!htb]
    \centering
    \caption{Exemplo de Quadro.\label{qua:quadro-exemplo1}}
    \begin{tabular}{|p{7cm}|p{7cm}|}
        \hline
        \textbf{BD Relacionais} & \textbf{BD Orientados a Objetos} \\
        \hline
        Os dados são passivos, ou seja, certas operações limitadas podem ser automaticamente acionadas quando os dados são usados. Os dados são ativos, ou seja, as solicitações fazem com que os objetos executem seus métodos. & Os processos que usam dados mudam constantemente. \\
        \hline
    \end{tabular}
    \fonte{\citeonline{Barbosa2004}}
\end{quadro}
 

A diferença entre quadro e tabela está no fato que um quadro é formado por linhas horizontais e verticais. Deve ser utilizado quando o conteúdo é majoritariamente não-numérico. O número do quadro e o título vem acima do quadro, e a fonte, deve vir abaixo. E Uma tabela é formada apenas por linhas verticais. Deve ser utilizada quando o conteúdo é majoritariamente numérico. O número da tabela e o título vem acima da tabela, e a fonte, deve vir abaixo, tal como no quadro.

\begin{table}[!htb]
    \centering
    \caption[Resultado dos testes]{Resultado dos testes.
    \label{tab:tabela-exemplo1}}
    \begin{tabular}{rrrrr}
        \toprule
            & Valores 1 & Valores 2 & Valores 3 & Valores 4 \\
        \midrule
            Caso 1 & 0,86 & 0,77 & 0,81 & 163 \\
            Caso 2 & 0,19 & 0,74 & 0,25 & 180 \\
            Caso 3 & 1,00 & 1,00 & 1,00 & 170 \\
        \bottomrule
    \end{tabular}
    \fonte{\citeonline{Barbosa2004}}
\end{table}


% EQUAÇÕES-----------------------------------------------------------------------
\section{EQUAÇÕES}
\label{sec:equacoes}

Esta seção trata de exemplos de como inserir a \autoref{eq:equacao-exemplo1} e a Eq. \ref{eq:equacao-exemplo2} no corpo do texto \footnote{Deve-se atentar ao fato de a formatação das equações ficar muito boa esteticamente.}. Observe que foram utilizadas duas formas distintas para referenciar as equações.

\begin{equation}
    X(s) = \int\limits_{t = -\infty}^{\infty} x(t) \, \text{e}^{-st} \, dt
    \label{eq:equacao-exemplo1}
\end{equation}

\begin{equation}
    F(u, v) = \sum_{m = 0}^{M - 1} \sum_{n = 0}^{N - 1} f(m, n) \exp \left[ -j 2 \pi \left( \frac{u m}{M} + \frac{v n}{N} \right) \right]
    \label{eq:equacao-exemplo2}
\end{equation}

% ALGORITMOS-----------------------------------------------------------------------
\section{ALGORITMOS}
\label{sec:algoritmos}

Esta seção trata de um exemplo de como inserir um algoritmo. Para inserção de algoritmos utiliza-se o pacote {\ttfamily algorithm2e} que já está devidamente configurado neste modelo.

Os algoritmos devem ser criados em arquivos separados para facilitar manutenção e armazenados no diretório de "/dados". O Algoritmo \ref{alg:algoritimo1} é um exemplo de algoritmo.  

\begin{algorithm}
    \label{alg:algoritimo1}
    \caption{Exemplo de Algoritmo}
    \KwIn{o número $n$ de vértices a remover, grafo original $G(V, E)$}
    \KwOut{grafo reduzido $G'(V,E)$}
    $removidos \leftarrow 0$ \\
    \While {removidos $<$ n } {
        $v \leftarrow$ Random$(1, ..., k) \in V$ \\
            \For {$u \in adjacentes(v)$} {
                remove aresta (u, v)\\
                $removidos \leftarrow removidos + 1$\\
            }
            \If {há  componentes desconectados} {
                remove os componentes desconectados\\
            }
        }
\end{algorithm}


% SOBRE AS LISTAS--------------------------------------------------------------------
\section{SOBRE AS LISTAS}
\label{sec:apSobreLista}

Para construir listas de "\textit{bullets}"{} ou listas enumeradas, inclusive listas aninhadas, é utilizado o pacote \verb|paralist|.

Exemplo de duas listas não numeradas aninhadas, utilizando o comando \verb|\itemize|. Observe a indentação, bem como a mudança automática do tipo de "\textit{bullet}"{} nas listas aninhadas.

\begin{itemize}
    \item item não numerado 1
    \item item não numerado 2
    \begin{itemize}
        \item subitem não numerado 1
        \item subitem não numerado 2
        \item subitem não numerado 3
    \end{itemize}
    \item item não numerado 3
\end{itemize}

Exemplo de duas listas numeradas aninhadas, utilizando o comando \verb|\enumerate|. Observe a numeração progressiva e indentação das listas aninhadas.

\begin{enumerate}
    \item item numerado 1
    \item item numerado 2
    \begin{enumerate}
        \item subitem numerado 1
        \item subitem numerado 2
        \item subitem numerado 3
    \end{enumerate}
    \item item numerado 3
\end{enumerate}

% SOBRE AS CITAÇÕES E CHAMADAS DE REFERÊNCAS----------------------------------------------
\section{SOBRE AS CITAÇÕES E CHAMADAS DE REFERÊNCAS}
\label{sec:apSobreCita}

Citações são trechos de texto ou informações obtidas de materiais consultadss quando da elaboração do trabalho. São utilizadas no texto com o propósito de esclarecer, completar e embasar as ideias do autor. Todas as publicações consultadas e utilizadas (por meio de citações) devem ser listadas, obrigatoriamente, nas referências bibliográficas, para preservar os direitos autorais. São classificadas em citações indiretas e diretas.

% CITAÇÕES INDIRETAS-----------------------------------------------------------------------
\subsection{CITAÇÕES INDIRETAS}
\label{subsec:citacoesLivres}

É a transcrição, com suas próprias palavras, das idéias de um autor, mantendo-se o sentido original. A citação indireta é a maneira que o pesquisador tem de ler, compreender e gerar conhecimento a partir do conhecimento de outros autores. Quanto à chamada da referência, ela pode ser feita de duas maneiras distintas, conforme o nome do(s) autor(es) façam parte do seu texto ou não. Exemplo de chamada fazendo parte do texto:\\
\\Enquanto \citeonline{Maturana2003} defendem uma epistemologia baseada na biologia. Para os autores, é necessário rever \ldots.\\

A chamada de referência foi feita com o comando \verb|\citeonline{chave}|, que produzirá a formatação correta.

A segunda forma de fazer uma chamada de referência deve ser utilizada quando se quer evitar uma interrupção na sequência do texto, o que poderia, eventualmente, prejudicar a leitura. Assim, a citação é feita e imediatamente após a obra referenciada deve ser colocada entre parênteses. Porém, neste caso específico, o nome do autor deve vir em caixa alta, seguido do ano da publicação. Exemplo de chamada não fazendo parte do texto:\\
\\Há defensores da epistemologia baseada na biologia que argumentam em favor da necessidade de \ldots \cite{Maturana2003}.\\

Nesse caso a chamada de referência deve ser feita com o comando \verb|\cite{chave}|, que produzirá a formatação correta.

% CITAÇÕES DIRETAS-----------------------------------------------------------------------
\subsection{CITAÇÕES DIRETAS}
\label{subsec:citacoesLiterais}

É a transcrição ou cópia de um parágrafo, de uma frase, de parte dela ou de uma expressão, usando exatamente as mesmas palavras adotadas pelo autor do trabalho consultado.

Quanto à chamada da referência, ela pode ser feita de qualquer das duas maneiras já mencionadas nas citações indiretas, conforme o nome do(s) autor(es) façam parte do texto ou não. Há duas maneiras distintas de se fazer uma citação direta, conforme o trecho citado seja longo ou curto.

Quando o trecho citado é longo (4 ou mais linhas) deve-se usar um parágrafo específico para a citação, na forma de um texto recuado (4 cm da margem esquerda), com tamanho de letra menor e espaçamento entrelinhas simples. Exemplo de citação longa:
\\\begin{citacao}
    Desse modo, opera-se uma ruptura decisiva entre a reflexividade filosófica, isto é a possibilidade do sujeito de pensar e de refletir, e a objetividade científica. Encontramo-nos num ponto em que o conhecimento científico está sem consciência. Sem consciência moral, sem consciência reflexiva e também subjetiva. Cada vez mais o desenvolvimento extraordinário do conhecimento científico vai tornar menos praticável a própria possibilidade de reflexão do sujeito sobre a sua pesquisa \cite[p.~28]{Silva2000}.
\end{citacao}

Para fazer a citação longa deve-se utilizar os seguintes comandos:
\begin{verbatim}
\begin{citacao}
<texto da citacao>
\end{citacao}
\end{verbatim}

No exemplo acima, para a chamada da referência o comando \verb|\cite[p.~28]{Silva2000}| foi utilizado, visto que os nomes dos autores não são parte do trecho citado. É necessário também indicar o número da página da obra citada que contém o trecho citado.

Quando o trecho citado é curto (3 ou menos linhas) ele deve inserido diretamente no texto entre aspas. Exemplos de citação curta:\\
\\A epistemologia baseada na biologia parte do princípio de que "assumo que não posso fazer referência a entidades independentes de mim para construir meu explicar" \cite[p.~35]{Maturana2003}.\\
\\A epistemologia baseada na biologia de \citeonline[p.~35]{Maturana2003} parte do princípio de que "assumo que não posso fazer referência a entidades independentes de mim para construir meu explicar".

% DETALHES SOBRE AS CHAMADAS DE REFERÊNCIAS---------------------------------------------------------
\section{DETALHES SOBRE AS CHAMADAS DE REFERÊNCIAS}
\label{sec:referUtilizadas}

Outros exemplos de comandos para as chamadas de referências e o resultado produzido por estes:\\
\\\citeonline{Maturana2003} \ \ \  \verb|\citeonline{Maturana2003}|\\
\citeonline{Barbosa2004} \ \ \   \verb|\citeonline{Barbosa2004}|\\
\cite[p.~28]{Silva2000} \ \ \  \verb|\cite[p.~28]{Silva2000}|\\
\citeonline[p.~33]{Silva2000} \ \ \   \verb|\citeonline[p.~33]{v}|\\
\cite[p.~35]{Maturana2003} \ \ \   \verb|\cite[p.~35]{Maturana2003}|\\
\citeonline[p.~35]{Maturana2003} \ \ \   \verb|\citeonline[p.~35]{Maturana2003}|\\
\cite{Barbosa2004,Maturana2003} \ \ \   \verb|\cite{Barbosa2004,Maturana2003}|\\

% SOBRE AS REFERÊNCIAS BIBLIOGRÁFICAS-------------------------------------------------------
\section{SOBRE AS REFERÊNCIAS BIBLIOGRÁFICAS}
\label{sec:apSobreRefer}

A bibliografia é feita no padrão \textsc{Bib}\TeX{}. As referências são colocadas em um arquivo separado. Neste template as referências são armazenadas no arquivo 

\begin{center}
``base-referencias.bib''.
\end{center}

Existem diversas categorias documentos e materiais componentes da bibliografia. A classe abn\TeX{} define as seguintes categorias (entradas):

\begin{verbatim}
@book
@inbook
@article
@phdthesis
@mastersthesis
@monography
@techreport
@manual
@proceedings
@inproceedings
@journalpart
@booklet
@patent
@unpublished
@misc
\end{verbatim}

Cada categoria (entrada) é formatada pelo pacote \citeonline{abnTeX22014d} de uma forma específica. Algumas entradas foram introduzidas especificamente para atender à norma \citeonline{NBR6023:2002}, são elas: \verb|@monography|, \verb|@journalpart|,\verb|@patent|. As demais entradas são padrão \textsc{Bib}\TeX{}. Para maiores detalhes, refira-se a \citeonline{abnTeX22014d}, \citeonline{abnTeX22014b}, \citeonline{abnTeX22014c}.

% NOTAS DE RODAPÉ--------------------------------------------------------------------------
\section{NOTAS DE RODAPÉ}
\label{sec:notasRodape}

As notas de rodapé pode ser classificadas em duas categorias: notas explicativas\footnote{é o tipo mais comum de notas que destacam, explicam e/ou complementam o que foi dito no corpo do texto, como esta nota de rodapé, por exemplo.} e notas de referências. A notas de referências, como o próprio nome ja indica, são utilizadas para colocar referências e/ou chamadas de referências sob certas condições.
           	% Orientações sobre o modelo Latex (comentar essa linha ao iniciar a escrita)
%---------pós texto-------------------------------------------------------------
\postextual
% REFERÊNCIAS------------------------------------------------------------------

% Carrega o arquivo "base-referencias.bib" e extrai automaticamente as referências citadas

\bibliography{./base-referencias}
\bibliographystyle{abntex2-alf} % Define o estilo ABNT para formatar a lista de referências
% OBSERVAÇÕES------------------------------------------------------------------
% Este arquivo não precisa ser alterado.
           	% Referências
% APÊNDICES--------------------------------------------------------------------

\begin{apendicesenv}
\partapendices

% Primeiro apêndice------------------------------------------------------------
\chapter{Nome do apêndice} % Edite para alterar o título deste apêndice
\label{chap:apendiceA}

Lembre-se que a diferença entre apêndice e anexo diz respeito à autoria do texto e/ou material ali colocado.

Caso o material ou texto suplementar ou complementar seja de sua autoria, então ele deverá ser colocado como um apêndice. Porém, caso a autoria seja de terceiros, então o material ou texto deverá ser colocado como anexo.

Caso seja conveniente, podem ser criados outros apêndices para o seu trabalho acadêmico. Basta recortar e colar este trecho neste mesmo documento. Lembre-se de alterar o "label"{} do apêndice.

Não é aconselhável colocar tudo que é complementar em um único apêndice. Organize os apêndices de modo que, em cada um deles, haja um único tipo de conteúdo. Isso facilita a leitura e compreensão para o leitor do trabalho.

% Novo apêndice----------------------------------------------------------------
\chapter{Nome do outro apêndice}
\label{chap:apendiceB}

conteúdo do novo apêndice

\end{apendicesenv}
                % Apêndices
% ANEXO------------------------------------------------------------------------

\begin{anexosenv}
\partanexos

% Primeiro anexo---------------------------------------------------------------
\chapter{Nome do anexo}     % edite para alterar o título deste anexo
\label{chap:anexoA}

Lembre-se que a diferença entre apêndice e anexo diz respeito à autoria do texto e/ou material ali colocado.

Caso o material ou texto suplementar ou complementar seja de sua autoria, então ele deverá ser colocado como um apêndice. Porém, caso a autoria seja de terceiros, então o material ou texto deverá ser colocado como anexo.

Caso seja conveniente, podem ser criados outros anexos para o seu trabalho acadêmico. Basta recortar e colar este trecho neste mesmo documento. Lembre-se de alterar o "label"{} do anexo.

Organize seus anexos de modo a que, em cada um deles, haja um único tipo de conteúdo. Isso facilita a leitura e compreensão para o leitor do trabalho. É para ele que você escreve.

% Novo anexo-------------------------------------------------------------------
\chapter{Nome do outro anexo}
\label{chap:anexoB}

conteúdo do outro anexo

\end{anexosenv}
               	% Anexos


\else % Se TCC2
% INSERE ELEMENTOS TEXTUAIS DO TCC2 (os elementos opcionais podem ser comentados)
% FOLHA DE APROVAÇÃO------------------------------------------------------------------

% Entre com os dados da banca
\newcommand{\hora}{10h40}
\newcommand{\dia}{30}
\newcommand{\mes}{fevereiro}
\newcommand{\ano}{\the\year}
\newcommand{\membroA}{Hommer Simpson\\ Instituição} % Membro A
\newcommand{\membroB}{Gordon Freeman\\ Instituição}   % Membro B
\newcommand{\coordenador}{Prof. Dr. Fábio Risental Coutinho}
\newcommand{\numerotcc}{70} % Pedir um número ao professor responsável pelo TCC

% Início da folha de aprovação (não alterar)
\begin{folhadeaprovacao}

    \noindent\begin{minipage}[t]{0.19\textwidth}
        \centering
        \raisebox{-0.8\height}{\includegraphics[width=0.7\textwidth]{dados/figuras/folha-aprovacao/brasao.eps}}
    \end{minipage}
    \begin{minipage}[t]{0.6\textwidth}
        \begin{center}
            Ministério da Educação\\
            \textbf{Universidade Tecnológica Federal do Paraná}\\
            Campus Toledo\\
            Coordenação do Curso de Engenharia Eletrônica
        \end{center}
    \end{minipage}
    \begin{minipage}[t]{0.19\textwidth}
        \centering
        \raisebox{-1.2\height}{\includegraphics[width=\textwidth]{dados/figuras/folha-aprovacao/logo_UTFPR_cor.eps}}
    \end{minipage}
    
    \vspace{0.5cm}
    
    \noindent\hrulefill
    
    \begin{center}
        TERMO DE APROVAÇÃO\\[5pt]
        Título do Trabalho de Conclusão de Curso \numero{\numerotcc{}}
    \end{center}
    %\vspace{1.0cm}
    \begin{center}
        \large\textbf{\imprimirtitulo}\\[10pt]
        por\\[10pt]
        \imprimirautor\\[15pt]
    \end{center}
        
    % Alterar o texto com a data e hora da defesa:
    \noindent Esse Trabalho de Conclusão de Curso foi apresentado às \textbf{\hora{} do dia \dia{} de \mes{} de \ano{}} como \textbf{requisito parcial} para a obtenção do título de \textbf{\imprimirtituloAcademico}. Após deliberação da Banca Examinadora, composta pelos professores abaixo assinados, o trabalho foi considerado \textbf{APROVADO}.
        
    % Assinaturas
    \begin{center}
        \assinatura{\membroA}
        \assinatura{\membroB}
        \assinatura{\imprimirorientador \\UTFPR}
        \assinatura{\coordenador \\ UTFPR}
    \end{center}
    
    % Finalização
 	\begin{center}
 	\small O termo de aprovação assinado encontra-se na coordenação do curso
	\vfill
	{\large\imprimirlocal, \dia{} de \mes{} de \ano{}}
	\end{center}
	
\end{folhadeaprovacao}
	

% DEDICATÓRIA------------------------------------------------------------------

\renewcommand{\dedicatorianame}{DEDICATÓRIA}

\begin{dedicatoria}

Altere este texto inserindo a dedicatória do seu trabalho. 

\end{dedicatoria}
          	% Dedicatória (opcional)
% AGRADECIMENTOS---------------------------------------------------------------

\begin{agradecimentos}[AGRADECIMENTOS]

Edite e coloque aqui os agradecimentos às pessoas e/ou instituições que contribuíram para a realização do trabalho.

É obrigatório o agradecimento às instituições de fomento à pesquisa que financiaram total ou parcialmente o trabalho, inclusive no que diz respeito à concessão de bolsas.

\end{agradecimentos}
        	% Agradecimentos (opcional)
% EPÍGRAFE---------------------------------------------------------------------

\renewcommand{\epigraphname}{EPÍGRAFE}

\begin{epigrafe}

\textit{Eu denomino meu campo de Gestão do Conhecimento, mas você não pode gerenciar conhecimento. Ninguém pode. O que pode fazer - o que a empresa pode fazer - é gerenciar o ambiente que otimize o conhecimento. (PRUSAK, Laurence, 1997).}

\end{epigrafe}

% OBSERVAÇÕES------------------------------------------------------------------
% Altere o texto para inserir a epígrafe do seu trabalho
              	% Epígrafe (opcional)
% RESUMO--------------------------------------------------------------------------------

\begin{resumo}[RESUMO]
\begin{SingleSpacing}

% Não altere esta seção do texto--------------------------------------------------------
%\imprimirautorcitacao. \imprimirtitulo. \imprimirdata. \pageref {LastPage} f. \imprimirprojeto\ – \imprimirprograma, \imprimirinstituicao. \imprimirlocal, \imprimirdata.\\
%---------------------------------------------------------------------------------------

O Resumo é um elemento obrigatório em tese, dissertação, monografia e TCC, constituído de uma seqüência de frases concisas e objetivas, fornecendo uma visão rápida e clara do conteúdo do estudo. O texto deverá conter no máximo 500 palavras e ser antecedido pela referência do estudo. Também, não deve conter citações. O resumo deve ser redigido em parágrafo único, espaçamento simples e seguido das palavras representativas do conteúdo do estudo, isto é, palavras-chave, em número de três a cinco, separadas entre si por ponto e finalizadas também por ponto. Usar o verbo na terceira pessoa do singular, com linguagem impessoal, bem como fazer uso, preferencialmente, da voz ativa. Texto contendo um único parágrafo.\\

\textbf{Palavras-chave}: Palavra. Segunda Palavra. Outra palavra.

\end{SingleSpacing}
\end{resumo}

% OBSERVAÇÕES---------------------------------------------------------------------------
% Altere o texto inserindo o Resumo do seu trabalho.
% Escolha de 3 a 5 palavras ou termos que descrevam bem o seu trabalho 
             		% Resumo em Português
% ABSTRACT--------------------------------------------------------------------------------

\begin{resumo}[ABSTRACT]
\begin{SingleSpacing}

% Não altere esta seção do texto--------------------------------------------------------
%\imprimirautorcitacao. \imprimirtitleabstract. \imprimirdata. \pageref {LastPage} f. \imprimirprojeto\ – \imprimirprograma, \imprimirinstituicao. \imprimirlocal, \imprimirdata.\\
%---------------------------------------------------------------------------------------

Elemento obrigatório em tese, dissertação, monografia e TCC. É a versão do resumo em português para o idioma de divulgação internacional. Deve ser antecedido pela referência do estudo. Deve aparecer em folha distinta do resumo em língua portuguesa e seguido das palavras representativas do conteúdo do estudo, isto é, das palavras-chave. Sugere-se a elaboração do resumo (Abstract) e das palavras-chave (Keywords) em inglês; para resumos em outras línguas, que não o inglês, consultar o departamento / curso de origem.\\

\textbf{Keywords}: Word. Second Word. Another word.

\end{SingleSpacing}
\end{resumo}

% OBSERVAÇÕES---------------------------------------------------------------------------
% Altere o texto inserindo o Abstract do seu trabalho.
% Escolha de 3 a 5 palavras ou termos que descrevam bem o seu trabalho 
             	% Resumo em Inglês
% Lista de Figuras----------------------------------------------------------------

\pdfbookmark[0]{\listfigurename}{lof}
\listoffigures*
\cleardoublepage

% OBSERVAÇÕES---------------------------------------------------------------------
% Este arquivo não precisa de ser alterado, pois a lista é gerada automaticamente.
   	% Lista de Figuras
% LISTA DE QUADROS----------------------------------------------------------------

\renewcommand{\listofquadrosname}{LISTA DE QUADROS}

\pdfbookmark[0]{\listofquadrosname}{loq}
\listofquadros*
\cleardoublepage

% OBSERVAÇÕES---------------------------------------------------------------------
% Este arquivo não necessita de ser editado. A lista é gerada automaticamente.
   	% Lista de Quadros (opcional)
% LISTA DE TABELAS-------------------------------------------------------------

\pdfbookmark[0]{\listtablename}{lot}
\listoftables*
\cleardoublepage

% OBSERVAÇÕES-------------------------------------------------------------------
% Este arquivo não precisa ser alterado, pois a lista é gerada automaticamente.
    % Lista de Tabelas
% LISTA DE ABREVIATURAS E SIGLAS----------------------------------------------------------

%\begin{siglas}
%    \item[ABNT] Associação Brasileira de Normas Técnicas
%    \item[DECOM] Departamento de Computação
%\end{siglas}

% OBSERVAÇÕES-----------------------------------------------------------------------------
% Altere a lista acima para definir os acrônimos e siglas utilizados neste trabalho


\listadesiglas     % Lista de Abreviaturas e Siglas
% LISTA DE SÍMBOLOS------------------------------------------------------------

%\begin{simbolos}
%    \item[$ \Gamma $] Letra grega Gama
%    \item[$ \lambda $] Comprimento de onda
%    \item[$ \in $] Pertence
%\end{simbolos}

% OBSERVAÇÕES-------------------------------------------------------------------
% Altere a lista acima para definir os símbolos utilizados no trabalho
\listadesimbolos   % Lista de Símbolos
% LISTA DE ALGORITMOS----------------------------------------------------------

\newcommand{\algoritmoname}{Algoritmo}
\renewcommand{\listalgorithmcfname}{LISTA DE ALGORITMOS}

\floatname{algocf}{\algoritmoname}
\newlistof{listofalgoritmos}{loa}{\listalgoritmoname}
\newlistentry{algocf}{loa}{0}

\counterwithout{algocf}{chapter}
\renewcommand{\cftalgocfname}{\algoritmoname\space}
\renewcommand*{\cftalgocfaftersnum}{\hfill--\hfill}

\pdfbookmark[0]{\listalgorithmcfname}{loa}
\listofalgorithms
\cleardoublepage

% OBSERVAÇÕES------------------------------------------------------------------
% Este arquivo não precisa ser alterado, pois a lista é gerada automaticamente.
 % Lista de Algoritmos (opcional)
% SUMÁRIO----------------------------------------------------------------------

\renewcommand{\contentsname}{SUMÁRIO}

\pdfbookmark[0]{\contentsname}{toc}
\tableofcontents*
\cleardoublepage

% OBSERVAÇÕES-------------------------------------------------------------------
% Este arquivo não precisa ser alterado, pois o sumário é gerado automaticamente.
               	% Sumário
%---------texto-----------------------------------------------------------------
\textual
% INTRODUÇÃO-------------------------------------------------------------------

\chapter{INTRODUÇÃO}
\label{chap:introducao}

Edite e coloque aqui o seu texto de introdução.

A Introdução é a parte inicial do texto, na qual devem constar o tema e a delimitação do assunto tratado, objetivos da pesquisa e outros elementos necessários para situar o tema do trabalho, tais como: justificativa, procedimentos metodológicos (classificação inicial), embasamento teórico (principais bases sintetizadas) e estrutura do trabalho, tratados de forma sucinta. Recursos utilizados e cronograma são incluídos quando necessário. Salienta-se que os procedimentos metodológicos e o embasamento teórico são tratados, posteriormente, em capítulos próprios e com a profundidade necessária ao trabalho de pesquisa.

\section{LEIA ESTA SEÇÃO ANTES DE COMEÇAR}
\label{sec:antesleiame}

Este documento é um \emph{template} \LaTeX{} que foi concebido, primariamente, para ser utilizado na elaboração de Trabalho de Conclusão de Curso \sigla[Trabalho de Conclusão de Curso]{TCC} em conformidade com as normas da Associação Brasileira de Normas técnicas \sigla[Associação Brasileira de Normas técnicas]{ABNT}.

Para a produção deste \emph{template} foi necessário adaptar o arquivo {\ttfamily abntex2.cls}. Assim, foi produzido o arquivo {\ttfamily utfpr-abntex2.cls} que define o \verb|documentclass| específico para a Universidade Tecnológica Federal do Paraná \sigla[Universidade Tecnológica Federal do Paraná]{UTFPR}.

Antes de começar a escrever o seu trabalho acadêmico utilizando este \emph{template}, é importante saber que há dois arquivos que você precisará editar para que a capa e a folha de rosto de seu trabalho sejam geradas automaticamente.
São eles os arquivos {\ttfamily utfpr-tcc.tex} e {\ttfamily capa.tex}, presentes no diretório raiz e {\ttfamily /comum}, respectivamente. 

No arquivo {\ttfamily utfpr-tcc.tex}, que contém toda a estrutura do documento, deverá ser informado a natureza do trabalho (TCC1 ou TCC2). Para isso, procure o seguinte trecho:

\begin{verbatim}	
		\newif\iftcctype
		\tcctypetrue	% Comente essa linha para TCC2		
\end{verbatim}

\noindent Comente a segunda linha para habilitar a compilação dos arquivos do TCC2, e não comente para os arquivos de TCC1. 

No arquivo {\ttfamily capa.tex} deverá ser informado nome do autor, título do trabalho, nome do orientador e outras informações necessárias.

Em sistemas Linux, a compilação para gerar um arquivo no formato pdf, incluindo corretamente as referências bibliográficas, deve ser realizada utilizando o comando \verb|makefile|, disponível na mesma pasta onde está o arquivo principal \verb|utfpr-tcc.tex|. Caso seja alterado o nome do arquivo \verb|utfpr-tcc.tex|, deverá ser alterado no arquivo \verb|makefile| também.

\section{ORGANIZAÇÃO DO TRABALHO}
\label{sec:organizacaoTrabalho}

Normalmente ao final da introdução é apresentada, em um ou dois parágrafos curtos, a organização do restante do trabalho acadêmico.
Deve-se dizer o quê será apresentado em cada um dos demais capítulos.

%\simb[texto]{x_i}  

%\sigla[exemplo]{EX}
               % Introdução
% OBJETIVOS------------------------------------------------------------------

\chapter{OBJETIVOS}
\label{chap:objetivos}

Os objetivos devem deixar claro o que se pretende alcançar com o trabalho proposto. 

\section{OBJETIVO GERAL}
\label{sec:obgeral}

Deve salientar o objetivo mais amplo do trabalho. Trata-se de uma ação e, portanto, deve ser construído sempre a partir de um verbo no infinitivo.

\section{OBJETIVOS ESPECÍFICOS}
\label{sec:obesp}

São as etapas a serem realizadas para que se atinja o objetivo geral. Sugere-se que na descrição dos objetivos específicos sejam utilizados verbos de ação:

\begin{itemize}
    \item Objetivos Específicos Descritivos: caracterizar, descrever, traçar, ...;
    \item Objetivos Específicos Exploratórios: conhecer, identificar, levantar, ...;   
    \item Objetivos Específicos Explicativos: analisar, avaliar, verificar, ...
\end{itemize}
                % Objetivos
% JUSTIFICATIVA------------------------------------------------------------------

\chapter{JUSTIFICATIVA}
\label{chap:justificativa}

Na justificativa deve estar claro:

\begin{itemize}
    \item Por que se pretende resolver o problema apresentado;
    \item Qual a relevância do trabalho dentro da área técnica;   
    \item Quais os possíveis ganhos que se pode agregar ao conhecimento técnico com a resolução deste problema;
		\item Quais os fatores motivadores para se dedicar a este problema específico.
\end{itemize}
            % Justificativa
% REVISÃO DE LITERATURA--------------------------------------------------------

\chapter{REVISÃO DA LITERATURA}
\label{chap:fundamentacaoTeorica}

Uma dica: é uma boa prática iniciar cada novo capítulo com um breve texto introdutório (tipicamente, dois ou três parágrafos) que deve deixar claro o quê será discutido no capítulo, bem como a organização do capítulo.

Este capítulo também pode ser chamado de ``REFERENCIAL TEÓRCIO''. Para mais informações sobre como estruturar um referencial teórico e evitar plágios, acesse a página: \href{https://guiadamonografia.com.br/como-fazer-referencial-teorico/}{O Guia da Monografia}


      % Referencial teorico
% MATERIAIS E MÉTODOS------------------------------------------------------------------

\chapter{MATERIAIS E MÉTODOS}
\label{chap:materiaisemetodos}

Este capítulo deve tratar de como se pretende realizar o trabalho/pesquisa:

\begin{itemize}
    \item Etapas;
    \item Procedimentos de estudos e coleta de dados;   
    \item Estratégias a serem utilizadas para análise de dados, realização de ensaios, etc.;
		\item Materiais que serão utilizados.
\end{itemize}        % Materiais e metodologia
% RESULTADOS ESPERADOS-------------------------------------------------------------------

\chapter{RESULTADOS ESPERADOS}
\label{resultadosesperados}

É nesse capítulo onde são colocados as expectativas do projeto. Para maiores informações acesse a página: \href{http://www.mundodamonografia.com.br/projeto-de-pesquisa-resultados/}{Mundo da Monografia}.

               % Resultados
% ORIENTAÇÕES GERAIS------------------------------------------------------------
\chapter{ORIENTAÇÕES GERAIS SOBRE O MODELO}
\label{chap:orientações}

Este capítulo trata de orientações gerais sobre como utilizar este modelo de TCC. Recomenda-se a leitura do texto compilado junto com o código \LaTeX{} deste capítulo.

% SOBRE AS LISTAS DE SÍMBOLOS E SIGLAS------------------------------------------
\section{SOBRE AS LISTAS DE SÍMBOLOS E SIGLAS} 
\label{sec:simbsiglas}

Símbolos são inseridos automaticamente na lista através do comando:

\begin{verbatim}
\simb[<significado>]{<símbolo>}
\end{verbatim}

%Exemplos de código são mostrados no \autoref{qua:simb}.

%{\input{./dados/quadros/simb.tex}}

Dessa forma, o código

\begin{verbatim}
\simb[Variação da tensão de saída]{\Delta V_x}
\end{verbatim}

\noindent resulta em ``$\Delta V_x$'' no texto e a adição deste termo na lista de símbolos.

Siglas e abreviaturas são inseridos automaticamente na lista através do comando:

\begin{verbatim}
\sigla[<significado>]{<sigla ou abreviatura>}
\end{verbatim}

Dessa forma, o código

\begin{verbatim}
\simb[Corrente Alternada]{CA}
\end{verbatim}

\noindent resulta em ``(CA)'' no texto e a adição deste termo na lista de Abreviaturas e Siglas. Note que a silga ou abreveatura fica entre parenteses.


% SOBRE AS REFERENCIAS CRUZADAS E ITENS NUMERADOS-------------------------------
\section{SOBRE AS REFERENCIAS CRUZADAS E ITENS NUMERADOS}
\label{sec:refs}

A numeração sequencial de figuras, tabelas, quadros, algoritmos e equações ocorre de modo automático. As ilustrações (figuras), tabelas, algoritmos e quadros serão indexadas automaticamente em suas respectivas listas. 

Referências cruzadas são obtidas através dos comandos \verb|\label{}| e \verb|\ref{}|. Sendo assim, não é necessário por exemplo, saber que o número de certo capítulo é \ref{chap:fundamentacaoTeorica} para colocar o seu número no texto. Outra forma que pode ser utilizada é esta: \autoref{chap:fundamentacaoTeorica}, facilitando a inserção, remoção e manejo de elementos numerados no texto sem a necessidade de renumerar todos esses elementos.

% FIGURAS-----------------------------------------------------------------------
\section{FIGURAS}
\label{sec:figuras}

Esta seção trata de um exemplo de como inserir uma figura. Note que a \autoref{fig:figura-exemplo1} aparece automaticamente na lista de figuras. Para saber mais sobre o uso de imagens no \LaTeX{} consulte literatura especializada \cite{Goossens2007}.

Os arquivos das figuras devem ser armazenados no diretório de "/dados", para facilitar a organização.

\begin{figure}[!htb]
    \centering
    \caption{Exemplo de Figura}
    \includegraphics[width=0.6\textwidth]{dados/figuras/orientacoes/blackmesalogo.eps}
    \fonte{\citeonline{valve1998}}
    \label{fig:figura-exemplo1}
\end{figure}

% QUADROS E TABELAS---------------------------------------------------------------
\section{QUADROS E TABELAS}
\label{sec:tabelas}

Esta seção trata de exemplos de como inserir o \autoref{qua:quadro-exemplo1} e a \autoref{tab:tabela-exemplo1}. Ambos aparecem automaticamente nas suas respectivas listas. Para saber mais informações sobre a construção de tabelas no \LaTeX{} consulte literatura especializada \cite{Mittelbach2004}.

Ambos os elementos (Quadros e Tabelas) devem ser criados em arquivos separados para facilitar manutenção e armazenados no diretório de "/dados".

\begin{quadro}[!htb]
    \centering
    \caption{Exemplo de Quadro.\label{qua:quadro-exemplo1}}
    \begin{tabular}{|p{7cm}|p{7cm}|}
        \hline
        \textbf{BD Relacionais} & \textbf{BD Orientados a Objetos} \\
        \hline
        Os dados são passivos, ou seja, certas operações limitadas podem ser automaticamente acionadas quando os dados são usados. Os dados são ativos, ou seja, as solicitações fazem com que os objetos executem seus métodos. & Os processos que usam dados mudam constantemente. \\
        \hline
    \end{tabular}
    \fonte{\citeonline{Barbosa2004}}
\end{quadro}
 

A diferença entre quadro e tabela está no fato que um quadro é formado por linhas horizontais e verticais. Deve ser utilizado quando o conteúdo é majoritariamente não-numérico. O número do quadro e o título vem acima do quadro, e a fonte, deve vir abaixo. E Uma tabela é formada apenas por linhas verticais. Deve ser utilizada quando o conteúdo é majoritariamente numérico. O número da tabela e o título vem acima da tabela, e a fonte, deve vir abaixo, tal como no quadro.

\begin{table}[!htb]
    \centering
    \caption[Resultado dos testes]{Resultado dos testes.
    \label{tab:tabela-exemplo1}}
    \begin{tabular}{rrrrr}
        \toprule
            & Valores 1 & Valores 2 & Valores 3 & Valores 4 \\
        \midrule
            Caso 1 & 0,86 & 0,77 & 0,81 & 163 \\
            Caso 2 & 0,19 & 0,74 & 0,25 & 180 \\
            Caso 3 & 1,00 & 1,00 & 1,00 & 170 \\
        \bottomrule
    \end{tabular}
    \fonte{\citeonline{Barbosa2004}}
\end{table}


% EQUAÇÕES-----------------------------------------------------------------------
\section{EQUAÇÕES}
\label{sec:equacoes}

Esta seção trata de exemplos de como inserir a \autoref{eq:equacao-exemplo1} e a Eq. \ref{eq:equacao-exemplo2} no corpo do texto \footnote{Deve-se atentar ao fato de a formatação das equações ficar muito boa esteticamente.}. Observe que foram utilizadas duas formas distintas para referenciar as equações.

\begin{equation}
    X(s) = \int\limits_{t = -\infty}^{\infty} x(t) \, \text{e}^{-st} \, dt
    \label{eq:equacao-exemplo1}
\end{equation}

\begin{equation}
    F(u, v) = \sum_{m = 0}^{M - 1} \sum_{n = 0}^{N - 1} f(m, n) \exp \left[ -j 2 \pi \left( \frac{u m}{M} + \frac{v n}{N} \right) \right]
    \label{eq:equacao-exemplo2}
\end{equation}

% ALGORITMOS-----------------------------------------------------------------------
\section{ALGORITMOS}
\label{sec:algoritmos}

Esta seção trata de um exemplo de como inserir um algoritmo. Para inserção de algoritmos utiliza-se o pacote {\ttfamily algorithm2e} que já está devidamente configurado neste modelo.

Os algoritmos devem ser criados em arquivos separados para facilitar manutenção e armazenados no diretório de "/dados". O Algoritmo \ref{alg:algoritimo1} é um exemplo de algoritmo.  

\begin{algorithm}
    \label{alg:algoritimo1}
    \caption{Exemplo de Algoritmo}
    \KwIn{o número $n$ de vértices a remover, grafo original $G(V, E)$}
    \KwOut{grafo reduzido $G'(V,E)$}
    $removidos \leftarrow 0$ \\
    \While {removidos $<$ n } {
        $v \leftarrow$ Random$(1, ..., k) \in V$ \\
            \For {$u \in adjacentes(v)$} {
                remove aresta (u, v)\\
                $removidos \leftarrow removidos + 1$\\
            }
            \If {há  componentes desconectados} {
                remove os componentes desconectados\\
            }
        }
\end{algorithm}


% SOBRE AS LISTAS--------------------------------------------------------------------
\section{SOBRE AS LISTAS}
\label{sec:apSobreLista}

Para construir listas de "\textit{bullets}"{} ou listas enumeradas, inclusive listas aninhadas, é utilizado o pacote \verb|paralist|.

Exemplo de duas listas não numeradas aninhadas, utilizando o comando \verb|\itemize|. Observe a indentação, bem como a mudança automática do tipo de "\textit{bullet}"{} nas listas aninhadas.

\begin{itemize}
    \item item não numerado 1
    \item item não numerado 2
    \begin{itemize}
        \item subitem não numerado 1
        \item subitem não numerado 2
        \item subitem não numerado 3
    \end{itemize}
    \item item não numerado 3
\end{itemize}

Exemplo de duas listas numeradas aninhadas, utilizando o comando \verb|\enumerate|. Observe a numeração progressiva e indentação das listas aninhadas.

\begin{enumerate}
    \item item numerado 1
    \item item numerado 2
    \begin{enumerate}
        \item subitem numerado 1
        \item subitem numerado 2
        \item subitem numerado 3
    \end{enumerate}
    \item item numerado 3
\end{enumerate}

% SOBRE AS CITAÇÕES E CHAMADAS DE REFERÊNCAS----------------------------------------------
\section{SOBRE AS CITAÇÕES E CHAMADAS DE REFERÊNCAS}
\label{sec:apSobreCita}

Citações são trechos de texto ou informações obtidas de materiais consultadss quando da elaboração do trabalho. São utilizadas no texto com o propósito de esclarecer, completar e embasar as ideias do autor. Todas as publicações consultadas e utilizadas (por meio de citações) devem ser listadas, obrigatoriamente, nas referências bibliográficas, para preservar os direitos autorais. São classificadas em citações indiretas e diretas.

% CITAÇÕES INDIRETAS-----------------------------------------------------------------------
\subsection{CITAÇÕES INDIRETAS}
\label{subsec:citacoesLivres}

É a transcrição, com suas próprias palavras, das idéias de um autor, mantendo-se o sentido original. A citação indireta é a maneira que o pesquisador tem de ler, compreender e gerar conhecimento a partir do conhecimento de outros autores. Quanto à chamada da referência, ela pode ser feita de duas maneiras distintas, conforme o nome do(s) autor(es) façam parte do seu texto ou não. Exemplo de chamada fazendo parte do texto:\\
\\Enquanto \citeonline{Maturana2003} defendem uma epistemologia baseada na biologia. Para os autores, é necessário rever \ldots.\\

A chamada de referência foi feita com o comando \verb|\citeonline{chave}|, que produzirá a formatação correta.

A segunda forma de fazer uma chamada de referência deve ser utilizada quando se quer evitar uma interrupção na sequência do texto, o que poderia, eventualmente, prejudicar a leitura. Assim, a citação é feita e imediatamente após a obra referenciada deve ser colocada entre parênteses. Porém, neste caso específico, o nome do autor deve vir em caixa alta, seguido do ano da publicação. Exemplo de chamada não fazendo parte do texto:\\
\\Há defensores da epistemologia baseada na biologia que argumentam em favor da necessidade de \ldots \cite{Maturana2003}.\\

Nesse caso a chamada de referência deve ser feita com o comando \verb|\cite{chave}|, que produzirá a formatação correta.

% CITAÇÕES DIRETAS-----------------------------------------------------------------------
\subsection{CITAÇÕES DIRETAS}
\label{subsec:citacoesLiterais}

É a transcrição ou cópia de um parágrafo, de uma frase, de parte dela ou de uma expressão, usando exatamente as mesmas palavras adotadas pelo autor do trabalho consultado.

Quanto à chamada da referência, ela pode ser feita de qualquer das duas maneiras já mencionadas nas citações indiretas, conforme o nome do(s) autor(es) façam parte do texto ou não. Há duas maneiras distintas de se fazer uma citação direta, conforme o trecho citado seja longo ou curto.

Quando o trecho citado é longo (4 ou mais linhas) deve-se usar um parágrafo específico para a citação, na forma de um texto recuado (4 cm da margem esquerda), com tamanho de letra menor e espaçamento entrelinhas simples. Exemplo de citação longa:
\\\begin{citacao}
    Desse modo, opera-se uma ruptura decisiva entre a reflexividade filosófica, isto é a possibilidade do sujeito de pensar e de refletir, e a objetividade científica. Encontramo-nos num ponto em que o conhecimento científico está sem consciência. Sem consciência moral, sem consciência reflexiva e também subjetiva. Cada vez mais o desenvolvimento extraordinário do conhecimento científico vai tornar menos praticável a própria possibilidade de reflexão do sujeito sobre a sua pesquisa \cite[p.~28]{Silva2000}.
\end{citacao}

Para fazer a citação longa deve-se utilizar os seguintes comandos:
\begin{verbatim}
\begin{citacao}
<texto da citacao>
\end{citacao}
\end{verbatim}

No exemplo acima, para a chamada da referência o comando \verb|\cite[p.~28]{Silva2000}| foi utilizado, visto que os nomes dos autores não são parte do trecho citado. É necessário também indicar o número da página da obra citada que contém o trecho citado.

Quando o trecho citado é curto (3 ou menos linhas) ele deve inserido diretamente no texto entre aspas. Exemplos de citação curta:\\
\\A epistemologia baseada na biologia parte do princípio de que "assumo que não posso fazer referência a entidades independentes de mim para construir meu explicar" \cite[p.~35]{Maturana2003}.\\
\\A epistemologia baseada na biologia de \citeonline[p.~35]{Maturana2003} parte do princípio de que "assumo que não posso fazer referência a entidades independentes de mim para construir meu explicar".

% DETALHES SOBRE AS CHAMADAS DE REFERÊNCIAS---------------------------------------------------------
\section{DETALHES SOBRE AS CHAMADAS DE REFERÊNCIAS}
\label{sec:referUtilizadas}

Outros exemplos de comandos para as chamadas de referências e o resultado produzido por estes:\\
\\\citeonline{Maturana2003} \ \ \  \verb|\citeonline{Maturana2003}|\\
\citeonline{Barbosa2004} \ \ \   \verb|\citeonline{Barbosa2004}|\\
\cite[p.~28]{Silva2000} \ \ \  \verb|\cite[p.~28]{Silva2000}|\\
\citeonline[p.~33]{Silva2000} \ \ \   \verb|\citeonline[p.~33]{v}|\\
\cite[p.~35]{Maturana2003} \ \ \   \verb|\cite[p.~35]{Maturana2003}|\\
\citeonline[p.~35]{Maturana2003} \ \ \   \verb|\citeonline[p.~35]{Maturana2003}|\\
\cite{Barbosa2004,Maturana2003} \ \ \   \verb|\cite{Barbosa2004,Maturana2003}|\\

% SOBRE AS REFERÊNCIAS BIBLIOGRÁFICAS-------------------------------------------------------
\section{SOBRE AS REFERÊNCIAS BIBLIOGRÁFICAS}
\label{sec:apSobreRefer}

A bibliografia é feita no padrão \textsc{Bib}\TeX{}. As referências são colocadas em um arquivo separado. Neste template as referências são armazenadas no arquivo 

\begin{center}
``base-referencias.bib''.
\end{center}

Existem diversas categorias documentos e materiais componentes da bibliografia. A classe abn\TeX{} define as seguintes categorias (entradas):

\begin{verbatim}
@book
@inbook
@article
@phdthesis
@mastersthesis
@monography
@techreport
@manual
@proceedings
@inproceedings
@journalpart
@booklet
@patent
@unpublished
@misc
\end{verbatim}

Cada categoria (entrada) é formatada pelo pacote \citeonline{abnTeX22014d} de uma forma específica. Algumas entradas foram introduzidas especificamente para atender à norma \citeonline{NBR6023:2002}, são elas: \verb|@monography|, \verb|@journalpart|,\verb|@patent|. As demais entradas são padrão \textsc{Bib}\TeX{}. Para maiores detalhes, refira-se a \citeonline{abnTeX22014d}, \citeonline{abnTeX22014b}, \citeonline{abnTeX22014c}.

% NOTAS DE RODAPÉ--------------------------------------------------------------------------
\section{NOTAS DE RODAPÉ}
\label{sec:notasRodape}

As notas de rodapé pode ser classificadas em duas categorias: notas explicativas\footnote{é o tipo mais comum de notas que destacam, explicam e/ou complementam o que foi dito no corpo do texto, como esta nota de rodapé, por exemplo.} e notas de referências. A notas de referências, como o próprio nome ja indica, são utilizadas para colocar referências e/ou chamadas de referências sob certas condições.
             % Orientações sobre o modelo Latex (comentar essa linha ao iniciar a escrita)
% CONCLUSÃO--------------------------------------------------------------------

\chapter{CONCLUSÃO}
\label{chap:conclusao}

Parte final do texto, na qual se apresentam as conclusões do trabalho acadêmico. É importante fazer uma análise crítica do trabalho, destacando os principais resultados e as contribuições do trabalho para a área de pesquisa.

\section{TRABALHOS FUTUROS}
\label{sec:trabalhosFuturos}

Também deve indicar, se possível e/ou conveniente, como o trabalho pode ser estendido ou aprimorado.

\section{CONSIDERAÇÕES FINAIS}
\label{sec:consideracoesFinais}

Encerramento do trabalho acadêmico.
                % Conclusão
%---------pós texto-------------------------------------------------------------
\postextual
% REFERÊNCIAS------------------------------------------------------------------

% Carrega o arquivo "base-referencias.bib" e extrai automaticamente as referências citadas

\bibliography{./base-referencias}
\bibliographystyle{abntex2-alf} % Define o estilo ABNT para formatar a lista de referências
% OBSERVAÇÕES------------------------------------------------------------------
% Este arquivo não precisa ser alterado.
           	% Referências
% APÊNDICES--------------------------------------------------------------------

\begin{apendicesenv}
\partapendices

% Primeiro apêndice------------------------------------------------------------
\chapter{Nome do apêndice} % Edite para alterar o título deste apêndice
\label{chap:apendiceA}

Lembre-se que a diferença entre apêndice e anexo diz respeito à autoria do texto e/ou material ali colocado.

Caso o material ou texto suplementar ou complementar seja de sua autoria, então ele deverá ser colocado como um apêndice. Porém, caso a autoria seja de terceiros, então o material ou texto deverá ser colocado como anexo.

Caso seja conveniente, podem ser criados outros apêndices para o seu trabalho acadêmico. Basta recortar e colar este trecho neste mesmo documento. Lembre-se de alterar o "label"{} do apêndice.

Não é aconselhável colocar tudo que é complementar em um único apêndice. Organize os apêndices de modo que, em cada um deles, haja um único tipo de conteúdo. Isso facilita a leitura e compreensão para o leitor do trabalho.

% Novo apêndice----------------------------------------------------------------
\chapter{Nome do outro apêndice}
\label{chap:apendiceB}

conteúdo do novo apêndice

\end{apendicesenv}
             	% Apêndices
% ANEXO------------------------------------------------------------------------

\begin{anexosenv}
\partanexos

% Primeiro anexo---------------------------------------------------------------
\chapter{Nome do anexo}     % edite para alterar o título deste anexo
\label{chap:anexoA}

Lembre-se que a diferença entre apêndice e anexo diz respeito à autoria do texto e/ou material ali colocado.

Caso o material ou texto suplementar ou complementar seja de sua autoria, então ele deverá ser colocado como um apêndice. Porém, caso a autoria seja de terceiros, então o material ou texto deverá ser colocado como anexo.

Caso seja conveniente, podem ser criados outros anexos para o seu trabalho acadêmico. Basta recortar e colar este trecho neste mesmo documento. Lembre-se de alterar o "label"{} do anexo.

Organize seus anexos de modo a que, em cada um deles, haja um único tipo de conteúdo. Isso facilita a leitura e compreensão para o leitor do trabalho. É para ele que você escreve.

% Novo anexo-------------------------------------------------------------------
\chapter{Nome do outro anexo}
\label{chap:anexoB}

conteúdo do outro anexo

\end{anexosenv}
               	% Anexos

\fi % fim do if

\end{document}
